\documentclass[a4paper,10pt]{article}
\usepackage[utf8]{inputenc}
\usepackage[T1]{fontenc}
\usepackage[french]{babel}
\usepackage[margin=1.5cm]{geometry}
\usepackage{xcolor}
\usepackage{tcolorbox}
\usepackage{enumitem}
\usepackage{multicol}
\usepackage{lmodern}
\usepackage{tabularx}
\usepackage{booktabs}
\usepackage{listings}

\renewcommand{\familydefault}{\sfdefault}
\pagestyle{empty}
\setlist{nosep, leftmargin=1.5em}

\definecolor{navy}{RGB}{0,51,102}
\definecolor{teal}{RGB}{0,120,120}
\definecolor{crimson}{RGB}{180,0,0}
\definecolor{green}{RGB}{0,120,50}
\definecolor{lightbg}{RGB}{245,248,252}
\definecolor{orange}{RGB}{200,100,0}
\definecolor{codebg}{RGB}{28,28,36}
\definecolor{codegreen}{RGB}{100,200,100}
\definecolor{codeorange}{RGB}{255,165,0}
\definecolor{codecyan}{RGB}{100,200,220}

\lstdefinestyle{algo}{
  backgroundcolor=\color{codebg},
  basicstyle=\ttfamily\small\color{white},
  keywordstyle=\color{codecyan}\bfseries,
  stringstyle=\color{codeorange},
  commentstyle=\color{gray}\itshape,
  frame=none, breaklines=true,
  morekeywords={ALGORITHME, Variables, DEBUT, FIN, Si, Sinon, FinSi, TantQue, FinTantQue, Pour, FinPour, Lire, Ecrire, Alors, Faire, jusqu, en},
}

\lstdefinestyle{csharp}{
  backgroundcolor=\color{codebg},
  basicstyle=\ttfamily\small\color{white},
  keywordstyle=\color{codecyan}\bfseries,
  stringstyle=\color{codeorange},
  commentstyle=\color{gray}\itshape,
  frame=none, breaklines=true,
  language=[Sharp]C,
}

\lstdefinestyle{python}{
  backgroundcolor=\color{codebg},
  basicstyle=\ttfamily\small\color{white},
  keywordstyle=\color{codecyan}\bfseries,
  stringstyle=\color{codeorange},
  commentstyle=\color{gray}\itshape,
  frame=none, breaklines=true,
  language=Python,
}

\tcbset{basebox/.style={
  colback=lightbg, colframe=navy, coltitle=white,
  fonttitle=\bfseries, arc=2mm, boxrule=0.5mm,
  left=3mm, right=3mm, top=2mm, bottom=2mm, after skip=6pt
}}

\newcommand{\key}[1]{\textbf{\textcolor{navy}{#1}}}
\newcommand{\warn}[1]{\textbf{\textcolor{crimson}{#1}}}

\begin{document}

\begin{center}
  \colorbox{teal}{\parbox{\dimexpr\textwidth-2\fboxsep}{
    \centering\vspace{3pt}
    {\LARGE\bfseries\textcolor{white}{PROGRAMMATION — Algorithmique, C\#, Python}}\\
    {\normalsize\textcolor{white}{BTS SIO — Maël Mignard}}
    \vspace{3pt}
  }}
\end{center}
\vspace{6pt}

\begin{multicols}{2}

% ============================================================
\begin{tcolorbox}[basebox, colframe=navy, title=1. Algorithmique — Structure générale]
\begin{lstlisting}[style=algo]
ALGORITHME NomAlgorithme
Variables
  age    : Numerique
  nom    : Chaine
  majeur : Booleen
DEBUT
  Ecrire "Entrez votre age : "
  Lire age
  majeur <- (age >= 18)
  Si majeur Alors
    Ecrire "Vous etes majeur."
  Sinon
    Ecrire "Vous etes mineur."
  FinSi
FIN
\end{lstlisting}
\key{Types} : Numérique, Chaîne, Booléen \\
\key{Opérateurs} : +, -, *, /, \% (modulo) \\
\key{Comparaison} : =, <>, <, >, <=, >= \\
\key{Logiques} : ET, OU, NON
\end{tcolorbox}

% ============================================================
\begin{tcolorbox}[basebox, colframe=teal, title=2. Algorithmique — Structures de contrôle]
\textbf{Condition :}
\begin{lstlisting}[style=algo]
Si condition Alors
  ...
SinonSi condition2 Alors
  ...
Sinon
  ...
FinSi
\end{lstlisting}
\vspace{4pt}
\textbf{Boucle TantQue (While) :}
\begin{lstlisting}[style=algo]
TantQue condition Faire
  ...
FinTantQue
\end{lstlisting}
\vspace{4pt}
\textbf{Boucle Pour (For) :}
\begin{lstlisting}[style=algo]
Pour i <- 1 jusqu'à 10 Faire
  Ecrire i
FinPour
\end{lstlisting}
\end{tcolorbox}

% ============================================================
\begin{tcolorbox}[basebox, colframe=orange, title=3. Algorigrammes — Symboles]
\begin{tabular}{ll}
\toprule
\textbf{Forme} & \textbf{Signification} \\
\midrule
Ovale          & Début / Fin \\
Rectangle      & Traitement / Action \\
Losange        & Condition / Test (Oui/Non) \\
Parallélogramme & Entrée (Lire) / Sortie (Ecrire) \\
Flèche         & Flux / Direction \\
\bottomrule
\end{tabular}
\vspace{4pt}\\
\key{Règle} : toujours commencer par Début et finir par Fin.\\
Un losange a \textbf{toujours 2 branches} : OUI et NON.
\end{tcolorbox}

% ============================================================
\begin{tcolorbox}[basebox, colframe=green, title=4. C\# — Bases du langage]
\begin{lstlisting}[style=csharp]
// Variables et types
int age = 20;
string nom = "Mael";
bool majeur = age >= 18;
double pi = 3.14;

// Affichage
Console.WriteLine($"Bonjour {nom}");

// Condition
if (age >= 18) {
    Console.WriteLine("Majeur");
} else {
    Console.WriteLine("Mineur");
}

// Boucle for
for (int i = 0; i < 5; i++) {
    Console.WriteLine(i);
}

// Boucle while
while (age < 65) {
    age++;
}
\end{lstlisting}
\end{tcolorbox}

% ============================================================
\begin{tcolorbox}[basebox, colframe=green, title=5. C\# — Tableaux et collections]
\begin{lstlisting}[style=csharp]
// Tableau statique
int[] notes = { 12, 15, 18, 10 };
Console.WriteLine(notes[0]); // 12

// foreach
foreach (int n in notes) {
    Console.WriteLine(n);
}

// Liste dynamique
List<string> noms = new List<string>();
noms.Add("Mael");
noms.Add("Camille");
Console.WriteLine(noms.Count); // 2

// Dictionnaire (clé/valeur)
Dictionary<string, int> ages = new();
ages["Mael"] = 20;
ages["Camille"] = 21;
\end{lstlisting}
\end{tcolorbox}

% ============================================================
\begin{tcolorbox}[basebox, colframe=navy, title=6. C\# — Programmation Orientée Objet]
\begin{lstlisting}[style=csharp]
public class Personne {
    // Propriétés
    public string Nom { get; set; }
    public int Age { get; set; }

    // Constructeur
    public Personne(string nom, int age) {
        Nom = nom;
        Age = age;
    }

    // Méthode
    public void SePresenter() {
        Console.WriteLine(
            $"Je suis {Nom}, j'ai {Age} ans.");
    }
}

// Instanciation
Personne p = new Personne("Mael", 20);
p.SePresenter();
\end{lstlisting}
\vspace{4pt}
\key{Héritage} : \texttt{class Etudiant : Personne} \\
\key{Interface} : contrat de méthodes (pas d'implémentation) \\
\key{Encapsulation} : \texttt{private}, \texttt{public}, \texttt{protected} \\
\key{Polymorphisme} : méthode redéfinie dans classe fille (\texttt{override})
\end{tcolorbox}

% ============================================================
\begin{tcolorbox}[basebox, colframe=teal, title=7. Python — Bases]
\begin{lstlisting}[style=python]
# Variables (typage dynamique)
age = 20
nom = "Mael"
majeur = age >= 18

# Affichage
print(f"Bonjour {nom}")

# Condition
if age >= 18:
    print("Majeur")
else:
    print("Mineur")

# Boucle for
for i in range(5):
    print(i)

# Boucle while
compteur = 0
while compteur < 10:
    compteur += 1

# Fonction
def aire_rectangle(l, h):
    return l * h

result = aire_rectangle(5, 3)  # 15
\end{lstlisting}
\end{tcolorbox}

% ============================================================
\begin{tcolorbox}[basebox, colframe=teal, title=8. Python — Listes et structures]
\begin{lstlisting}[style=python]
# Liste
notes = [12, 15, 18, 10]
notes.append(14)
print(notes[0])  # 12
print(len(notes))  # 5

# Dictionnaire
personne = {
    "nom": "Mael",
    "age": 20
}
print(personne["nom"])

# Compréhension de liste
carres = [x**2 for x in range(10)]

# Tri
notes.sort()
notes.sort(reverse=True)
\end{lstlisting}
\end{tcolorbox}

% ============================================================
\begin{tcolorbox}[basebox, colframe=orange, title=9. Comparaison Algorithme / C\# / Python]
\begin{tabularx}{\columnwidth}{lXX}
\toprule
\textbf{Concept} & \textbf{C\#} & \textbf{Python} \\
\midrule
Afficher & \texttt{Console.Write} & \texttt{print()} \\
Lire & \texttt{Console.ReadLine} & \texttt{input()} \\
Si/Sinon & \texttt{if / else \{\}} & \texttt{if / else:} \\
Boucle for & \texttt{for(int i=0;..)} & \texttt{for i in range:} \\
Fonction & \texttt{public int Fn()} & \texttt{def fn():} \\
Commentaire & \texttt{// texte} & \texttt{\# texte} \\
Tableau & \texttt{int[] tab} & \texttt{liste = []} \\
\bottomrule
\end{tabularx}
\end{tcolorbox}

% ============================================================
\begin{tcolorbox}[basebox, colframe=crimson, title=10. Concepts fondamentaux POO]
\begin{tabular}{ll}
\toprule
\textbf{Concept} & \textbf{Description} \\
\midrule
Classe & Modèle/plan d'un objet \\
Objet & Instance d'une classe \\
Attribut & Variable d'une classe \\
Méthode & Fonction d'une classe \\
Constructeur & Méthode d'initialisation \\
Héritage & Une classe réutilise une autre \\
Encapsulation & Masquer les détails internes \\
Polymorphisme & Même méthode, comportements différents \\
Abstraction & Simplifier la complexité \\
\bottomrule
\end{tabular}
\end{tcolorbox}

\end{multicols}

\vspace{4pt}
\begin{tcolorbox}[basebox, colframe=teal, title=Points clés à retenir pour l'examen]
\begin{multicols}{2}
\begin{itemize}
  \item Un algorithme = structure ALGORITHME / Variables / DEBUT / FIN
  \item Losange dans algorigramme = condition (2 sorties OUI/NON)
  \item \warn{TantQue} vérifie AVANT ; \warn{Répéter...Jusqu'à} vérifie APRÈS
  \item En C\# : types statiques (int, string, bool), accolades \{\}
  \item En Python : typage dynamique, indentation obligatoire
  \item POO : classe = plan, objet = instance
  \item Héritage C\# : \texttt{class Etudiant : Personne}
  \item Console.WriteLine = afficher + saut de ligne
  \item \texttt{range(0, 10)} = de 0 à 9 inclus
  \item Un constructeur s'appelle avec \texttt{new NomClasse()}
\end{itemize}
\end{multicols}
\end{tcolorbox}

\end{document}
